\section{A Modular Serial Concatenation Framework}
\label{sec:sc-framework}

% Uses notation/macros from the current Section 3:
% \Cone,\Ctwo,\Gone,\Gtwo,\enum.


We study depth-2 serially concatenated convolutional codes
\[
\Csc \;:=\; \Ctwo \circ \Pi \circ \Cone \;\subseteq\; \{0,1\}^n,
\]
where \(\Pi\) is a uniform random interleaver (uniform random permutation) on \([n]\).
Equivalently, writing generator matrices \(\Gone\in\{0,1\}^{k\times n}\), \(\Pi\in\{0,1\}^{n\times n}\),
and \(\Gtwo\in\{0,1\}^{n\times n}\), a codeword can be computed as
\[
y \;=\; x\cdot \Gone \cdot \Pi \cdot \Gtwo \in \{0,1\}^n
\qquad\text{for a message }x\in\{0,1\}^k.
\]
(Here \(\Ctwo\) is rate-1, hence represented by an \(n\times n\) matrix. In later sections \(\Ctwo\) will be a
recursive convolutional encoder, possibly drawn from a random inner ensemble.)

We follow the standard convention: \(\Cone\) is the outer code and \(\Ctwo\) is the inner code.

\paragraph{Goal of this section.}
This section factors the minimum-distance analysis into:
(i) an outer interface (logarithmic distance + spectral growth), and
(ii) an inner interface (a contraction bound for small output weight on a random interleaved input).
Once these interfaces are established, the linear-distance conclusion follows by a one-page \stan{check it is one page} first-moment argument.
All later instantiations (different outers and different inners) plug into this same framework.

% ============================================================
\subsection{Outer interface: distance and spectrum}
\label{sec:outer-interface}

Let \(\Cone=\Cone(n)\subseteq\{0,1\}^n\) be a (possibly random) linear code of rate \(R_1\).
For each weight \(w\in\{0,1,\dots,n\}\), define the (expected) outer weight spectrum
\[
\enum_{\Cone,w} \;:=\; \E\Bigl[\bigl|\{c\in\Cone:\ \wt(c)=w\}\bigr|\Bigr],
\]
where the expectation is over the randomness of \(\Cone\) (and is omitted if \(\Cone\) is deterministic).

We will use the following two outer properties throughout the paper.




\begin{assumption}[Outer interface]
\label{ass:outer-interface}
Fix an outer-code ensemble \(\Cone=\Cone(n)\).
There exist constants \(\alpha>0\), \(c\ge 0\), and \(\beta>1\), depending on the ensemble but not on \(n\),
such that for all sufficiently large \(n\):
\begin{enumerate}
\item \textbf{(O1) Logarithmic distance:} \(d_{\min}(\Cone)\ge \alpha\log n\) with probability \(1-o(1)\).
\item \textbf{(O2) Spectral growth:} \(\enum_{\Cone,w}\le n^c\,\beta^w\) for all \(w\in\{1,\dots,n\}\).
\end{enumerate}
\end{assumption}




\noindent
Later sections instantiate \(\Cone\) by two different families:
a log-memory dense terminated convolutional outer, and a sparse-memory (local-expander / banded LDPC) outer.

% ============================================================
\subsection{Inner interface: contraction under a random interleaver}
\label{sec:inner-interface}

The role of the random interleaver is to randomize locations so that, conditioned on an outer weight \(w\),
the input to the inner code is uniform over the Hamming slice of weight \(w\).

\begin{lemma}[Interleaver makes a fixed weight-\(w\) word uniform on the slice]
\label{lem:interleaver-uniform-slice}
Fix any \(c\in\{0,1\}^n\) with \(\wt(c)=w\). If \(\Pi\) is a uniform random permutation of \([n]\),
then \(\Pi(c)\) is uniform over \(\{u\in\{0,1\}^n:\ \wt(u)=w\}\).
\end{lemma}

\noindent
\emph{Proof.}
Every weight-\(w\) vector is obtained from \(c\) by exactly \(w!\,(n-w)!\) permutations, and \(\Pi\) is uniform on \(S_n\). \(\square\)

\paragraph{Slice-to-tail probability (the object we need from the inner).}
Let \(U\) be uniform over \(\{0,1\}^n\) conditioned on \(\wt(U)=w\).
Define, for \(\delta\in(0,1)\),
\begin{equation}
\label{eq:pw-def}
p_w(\delta)
\;:=\;
\Pr\bigl[\wt(\Ctwo(U))\le \delta n\ \big|\ \wt(U)=w\bigr].
\end{equation}
If \(\Ctwo\) itself is drawn from a random inner ensemble (e.g. dense time-varying recursion \ref{...}, sparse recursion with refresh \ref{...}),
then the probability above is taken over both the inner randomness and \(U\).
(Equivalently, one may interpret \(p_w(\delta)\) as the average slice-to-tail probability under the inner ensemble.)

\stan{in 3.1 we say how C1 is instantiated (we give the 2 options). We should say here how C2 is?}

\paragraph{Optional: an inner input--output enumerator viewpoint.}
When \(\Ctwo\) is deterministic and linear, one can also define the (slice) enumerator
\[
\enum^{\Ctwo}_{w,h}
\;:=\;
\bigl|\{u\in\{0,1\}^n:\ \wt(u)=w,\ \wt(\Ctwo(u))=h\}\bigr|,
\]
so that
\(
\Pr[\wt(\Ctwo(U))=h\mid \wt(U)=w] = \enum^{\Ctwo}_{w,h}/\binom{n}{w}
\)
and \(p_w(\delta)=\sum_{h\le \delta n}\enum^{\Ctwo}_{w,h}/\binom{n}{w}\).
Our accumulator analysis will proceed in this style; the other inners are analyzed directly via martingales / mixing.

\begin{assumption}[Inner interface]
\label{ass:inner-interface}
Fix a target relative distance $\delta\in(0,1)$.
There exist constants $\rho=\rho(\delta)\in(0,1)$ and $a\ge 0$ such that for all sufficiently large $n$ and all
\[
w \ \ge\ d_{\min}(\Cone),
\]
if $U$ is uniform over $\{0,1\}^n$ conditioned on $\wt(U)=w$, then the inner encoder $\Ctwo$ satisfies
\begin{equation}
\label{eq:inner-interface-pw}
p_w(\delta)
\;:=\;
\PP\big[\wt(\Ctwo(U))\le \delta n \ \big|\ \wt(U)=w\big]
\;\le\;
n^{a}\,\rho^{\,w}.
\end{equation}
\end{assumption}


\noindent
This is the formal version of the ``interleaver gain'' / small-output contraction phenomenon:
a random interleaved weight-\(w\) input produces an output below \(\delta n\) with probability exponentially small in \(w\),
up to a polynomial prefactor.

% ============================================================
\subsection{The SC framework theorem (first-moment reduction)}
\label{sec:framework-theorem}

\begin{theorem}[SC framework theorem]
\label{thm:sc-framework}
Fix \(\delta\in(0,1)\). Let \(\Cone\subseteq\{0,1\}^n\) be any fixed linear outer code and let \(\Ctwo\) be any (possibly random)
rate-1 linear inner encoder independent of the interleaver \(\Pi\).
Let \(\Csc=\Ctwo\circ\Pi\circ\Cone\).

Define the random variable counting low-weight nonzero concatenated codewords:
\[
X_\delta
\;:=\;
\bigl|\{c\in\Cone\setminus\{0\}:\ \wt(\Ctwo(\Pi(c)))\le \delta n\}\bigr|.
\]
Then
\begin{equation}
\label{eq:first-moment-master}
\Pr\bigl[d_{\min}(\Csc)\le \delta n\bigr]
\;\le\;
\EE[X_\delta]
\;=\;
\sum_{w=d_{\min}(\Cone)}^{n} \enum_{\Cone,w}\cdot p_w(\delta),
\end{equation}
where \(p_w(\delta)\) is defined in \eqref{eq:pw-def}.
\end{theorem}

\begin{proof}
The event \(d_{\min}(\Csc)\le \delta n\) occurs iff there exists a nonzero
\(c\in\Cone\) such that \(\wt(\Ctwo(\Pi(c)))\le \delta n\), i.e. iff \(X_\delta\ge 1\).
Thus by Markov's inequality,
\[
\Pr[d_{\min}(\Csc)\le \delta n]
\;=\;
\Pr[X_\delta\ge 1]
\;\le\;
\EE[X_\delta].
\]
To compute \(\EE[X_\delta]\), group outer codewords by weight \(w\).
For each fixed \(c\in\Cone\) with \(\wt(c)=w\), Lemma~\ref{lem:interleaver-uniform-slice} implies that \(\Pi(c)\)
is uniform over the Hamming slice of weight \(w\), hence
\[
\Pr[\wt(\Ctwo(\Pi(c)))\le \delta n \mid \wt(c)=w] \;=\; p_w(\delta),
\]
which depends only on \(w\). Summing over all codewords of each weight yields
\(\EE[X_\delta]=\sum_w \enum_{\Cone,w}p_w(\delta)\), and \eqref{eq:first-moment-master} follows. \(\square\)
\end{proof}

\paragraph{Why this avoids a union bound over messages.}
Equation \eqref{eq:first-moment-master} is a first-moment bound over codewords of \(\Cone\),
organized by the outer spectrum \(\enum_{\Cone,w}\). This is the mechanism that prevents an explicit union bound over
\(2^k\) messages in later minimum-distance arguments.

============================================================
\subsection{A convenient corollary under the outer/inner interfaces}
\label{sec:framework-corollary}

\begin{corollary}[Linear distance from outer spectrum + inner contraction]
\label{cor:linear-dist-from-interfaces}
Assume the outer interface (Assumption~\ref{ass:outer-interface}) and the inner interface
(Assumption~\ref{ass:inner-interface}) hold with parameters \((\alpha,c,\beta)\) and \((a,\rho)\), respectively.
If \(\beta\rho<1\) and
\begin{equation}
\label{eq:poly-beating-condition}
\alpha\cdot \log\Big(\frac{1}{\beta\rho}\Big) \;>\; a+c+2,
\end{equation}
then
\[
\Pr\bigl[d_{\min}(\Csc)\le \delta n\bigr] \;=\; o(1)
\qquad\text{as }n\to\infty.
\]
In particular, \(d_{\min}(\Csc)=\Omega(n)\) with high probability (over the interleaver and any inner randomness).
\end{corollary}

\begin{proof}
By Theorem~\ref{thm:sc-framework},
\[
\EE[X_\delta]
\;=\;
\sum_{w\ge d_{\min}(\Cone)} \enum_{\Cone,w}\cdot p_w(\delta).
\]
Using Assumption~\ref{ass:outer-interface} (spectrum bound) and Assumption~\ref{ass:inner-interface}
(which applies for all \(w\ge d_{\min}(\Cone)\)),
\[
\EE[X_\delta]
\;\le\;
\sum_{w\ge d_{\min}(\Cone)} (n^c\beta^w)\cdot(n^a\rho^w)
\;=\;
n^{a+c}\sum_{w\ge d_{\min}(\Cone)} (\beta\rho)^w.
\]
Since \(\beta\rho<1\), the geometric series is bounded by
\[
\sum_{w\ge d_{\min}(\Cone)}(\beta\rho)^w
\;\le\;
\frac{(\beta\rho)^{d_{\min}(\Cone)}}{1-\beta\rho}
\;\le\;
O(1)\cdot (\beta\rho)^{\alpha\log n}
\;=\;
O(1)\cdot n^{-\alpha\log(1/(\beta\rho))},
\]
using \(d_{\min}(\Cone)\ge \alpha\log n\).
Thus
\[
\EE[X_\delta]
\;\le\;
O(1)\cdot n^{a+c-\alpha\log(1/(\beta\rho))},
\]
which is \(o(1)\) exactly under \eqref{eq:poly-beating-condition}. The claim follows by Markov. \(\square\)
\end{proof}






\paragraph{How the rest of the paper uses the framework.}
From this point on, the paper decomposes into independent tasks:
\begin{itemize}
\item \textbf{Outer instantiation(s):} construct outer ensembles satisfying Assumption~\ref{ass:outer-interface}
(logarithmic distance and spectral growth), either with log memory (dense) or sparse memory (LDPC/local-expander).
\item \textbf{Inner instantiation(s):} prove contraction bounds of the form \eqref{eq:inner-contraction-template}
for each candidate inner (accumulator; dense time-varying recursion; sparse recursion with periodic refresh).
\end{itemize}
Each combination of a valid outer and a valid inner yields a linear-distance SC ensemble via
Corollary~\ref{cor:linear-dist-from-interfaces}.
